\subsection{The weighted arithmetic-path budget}
\label{sec:pathwise-guard}

The precision analysis uses the length of an arithmetic dependency path
through the framed complex network. We first measure this path in
directional rank units, before grouping some units into bulk children.

\paragraph{A weighted path bound.}
Use the retained triple-based complex motif on $h$ points, with
$m=h^3$. This includes the paired complex producer and its shared
first- and third-stage auxiliary banks. Regard every scalar gate as a
map from saved inputs to its outputs. All of its inputs and outputs
have one common subspace label. A directed dependency path can therefore
enter a gate on one physical wire and leave on another without changing
label dimension at the gate itself.

For a network edge $a$, let $U_a^-$ and $U_a^+$ be its labels and write
$r_a=|\dim U_a^+-\dim U_a^-|$. The edge implementation has exactly
$r_a$ recursive calls, each implementing one directional $C$ tensor
on $e/m$ selected axes. Binary changes of coordinates before and after
these calls permute coefficients and add no arithmetic depth. An inverse
direction is implemented using one forward child and its sign and
fourth-root wrappers, so it has the same recursive-call count.

If $P$ is a directed scalar dependency path, its common gate labels give
the exact identity
\[
 \sum_{a\in P}r_a
 =\dim U_{\rm end}-\dim U_{\rm start}
   +2\sum_{\substack{a\in P\\\dim U_a^+<\dim U_a^-}}
       (\dim U_a^--\dim U_a^+).
\]
The endpoint difference is at most $m$.

Every side-producer edge is increasing. This follows from nested source
labels for the nonconstant middle mixer, complemented labels in its
reverse-direction counterpart, and constant early and late frames.
The data-wire and interstage edges are increasing as in the retained
motif. The stage-sharing joins are also increasing. The only decreasing
edges are the retained central returns, each of dimension loss $h$.

There are three tensor stages. Within a stage, the $v^2$ invocations
use disjoint data and auxiliary wires, and each invocation has a single
central-return frontier. Equivalently, its point and total central wires
all return between the same two blocks of the transparent schedule.
A directed path crosses that frontier at most once. Gates acting on
different invocations do not connect their wires. Thus a path encounters
at most one central return per stage, even though the machine executes
all invocations sequentially on its fixed tapes. Sharing stage-1 and
stage-3 banks only links their existing stage boundaries and creates no
additional return. Consequently
\begin{equation}\label{eq:pathwise-child-bound}
 \sum_{a\in P}r_a\le q:=m+6h.
\end{equation}
The bound also applies to a path ending at an intermediate gate or inside
an edge implementation: complete its current edge, and then any suffix
to a terminal vertex. All weights are nonnegative.

\paragraph{From a scalar path to a coefficient path.}
A coefficient dependency path in the actual tape algorithm projects to
a directed path of the scalar network: inside a scalar gate it follows
one of the saved operands that contributes to the chosen output; inside
an edge it stays on the edge's role stream. A child can permute or combine
addresses within that role, but cannot move a value into a different
parent role. The row split, parking, padding and reassembly operations
do not create arithmetic dependencies between distinct parent roles.
This projection justifies using \eqref{eq:pathwise-child-bound} for the
implemented recursion, rather than merely for a drawn finite circuit.
